\section{Cambio de base en las categorías sobre~$\mathcal{E}$}
\label{VI.3}

Como en~$\Cat$ existen los límites proyectivos (respecto de categorías~$\cal{I}$
que son elementos de~$\Univ$), lo mismo ocurre en~$\Cat_{/\cal{E}}$; en
particular, existen allí los productos cartesianos, que se interpretan como
productos fibrados en~$\Cat$. De acuerdo con las notaciones generales, si
$\cal{F}$ y $\cal{G}$ son categorías sobre~$\cal{E}$, se denota por
$$
\cal{F}\times_{\cal{E}}\cal{G}
$$
\label{indnot:fk}\oldindexnot{$\cal{F}\times_{\cal{E}}\cal{G}$|hyperpage}%
su producto en~$\Cat_{/\cal{E}}$, es decir, su producto fibrado sobre~$\cal{E}$
en~$\Cat$,
\marginpar{153}
considerado como categoría sobre~$\cal{E}$. Así,
$\cal{F}\times_{\cal{E}}\cal{G}$ está provisto de dos $\cal{E}$-funtores
$\pr_1$ y $\pr_2$, que definen, para toda categoría $\cal{H}$ sobre~$\cal{E}$,
una biyección
$$
\Hom_{\cal{E}}(\cal{H},\cal{F}\times_{\cal{E}}\cal{G})\isomto
\Hom_{\cal{E}}(\cal{H},\cal{F})\times\Hom_{\cal{E}}(\cal{H},\cal{G}).
$$
Esta biyección procede, además, de un isomorfismo de categorías
$$
\SheafHom_{\cal{E}/-}(\cal{H},\cal{F}\times_{\cal{E}}
\cal{G})\isomto\SheafHom_{\cal{E}/-}(\cal{H},\cal{F})\times
\SheafHom_{\cal{E}/-}(\cal{H},\cal{G})
$$
al tomar los conjuntos de objetos de ambos miembros, donde el funtor escrito es
el que tiene por componentes los funtores $h\mto\pr_1\circ h$ y
$h\mto\pr_2\circ h$ del primer miembro en los dos factores del segundo. Se
deja al lector comprobar que se obtiene efectivamente un isomorfismo (pues el
hecho análogo es válido, más generalmente, siempre que se tenga un límite
proyectivo de categorías, y no solo en el caso de un producto fibrado).

Recordemos, además, que (como se dijo en el~\No \Ref{VI.1}):
\begin{eqnarray*}
\Ob(\cal{F}\times_\cal{E}\cal{G})&=&
\Ob(\cal{F})\times_{\Ob(\cal{E})}\Ob(\cal{G})\\
\Fl(\cal{F}\times_\cal{E}\cal{G})&=&
\Fl(\cal{F})\times_{\Fl(\cal{E})}\Fl(\cal{G})\quoi,
\end{eqnarray*}
y la composición de las flechas se realiza componente a componente.

En lo que sigue, consideraremos un funtor
$$
\lambda\colon\cal{E}'\to\cal{E}
$$
y, para toda categoría $\cal{F}$ sobre~$\cal{E}$, consideraremos
$\cal{F}\times_\cal{E}\cal{E}'$ como categoría sobre~$\cal{E}'$ por medio
de~$\pr_2$; en otras palabras, interpretaremos la operación «producto fibrado»
como una operación de \emph{«cambio de base»}, y el funtor
$\lambda\colon\cal{E}'\to\cal{E}$ recibe el nombre de \emph{«funtor de cambio
de base»}.
\index{cambio de base (funtor)|hyperpage}%
\index{funtor de cambio de base|hyperpage}%
De acuerdo con los hechos generales bien conocidos, se obtiene así un funtor,
llamado
\marginpar{154}
\emph{funtor de cambio de base} para~$\lambda$:
$$
\lambda^*\colon\Cat_{/\cal{E}}\to\Cat_{/\cal{E}'}\quoi,
$$
\label{indnot:fl}\oldindexnot{$\lambda^*\colon\Cat_{/\cal{E}}\to\Cat_{/\cal{E}'}$|hyperpage}%
(adjunto al funtor de «restricción de la base» que, a toda categoría~$\cal{F}'$
sobre~$\cal{E}'$ con funtor de proyección~$p'$, asocia~$\cal{F}'$, considerada
como categoría sobre~$\cal{E}$ mediante el funtor~$p=\lambda p'$). Como es bien
sabido en el caso general de un funtor de cambio de base en una categoría, el
funtor de cambio de base «conmuta con los límites proyectivos» y, en particular,
«transforma» productos fibrados sobre~$\cal{E}$ en productos fibrados
sobre~$\cal{E}'$.

Sean $\cal{F}$ y $\cal{G}$ dos categorías sobre~$\cal{E}$; vamos a definir un
\emph{isomorfismo canónico}:
\begin{multline*}
\label{eq:VI.3.i}
\tag{i}
\SheafHom_{\cal{E}'/-}(\cal{F}',\cal{G}')\isomto
\SheafHom_{\cal{E}/-}(\cal{F}\times_{\cal{E}}\cal{E}',\cal{G})\\
\text{donde }\cal{F}'=\cal{F}\times_{\cal{E}}\cal{E}'\text{, }
\cal{G}'=\cal{G}\times_{\cal{E}}\cal{E}'.
\end{multline*}
Para ello, consideremos el funtor
$$
\pr_1\colon\cal{G}'=\cal{G}\times_{\cal{E}}\cal{E}'\to\cal{G}\quoi,
$$
y definamos~\eqref{eq:VI.3.i} por
$$
F\mto\pr_1\circ F\quoi,
$$
que, a priori, designa un funtor
\begin{equation*}
\label{eq:VI.3.ii}
\tag{ii} {\SheafHom(\cal{F}',\cal{G}')\to\SheafHom(\cal{F}',\cal{G})}
\end{equation*}
Solo hay que comprobar, por tanto, que este último induce un funtor en las
subcategorías~\eqref{eq:VI.3.i} y que dicho funtor es un isomorfismo. Que
\eqref{eq:VI.3.ii} induce una biyección
$$
\Hom_{\cal{E}'/-}(\cal{F}',\cal{G}')\isomto
\Hom_{\cal{E}/-}(\cal{F}\times_{\cal{E}}\cal{E}',\cal{G})
$$
es la propiedad característica del funtor de cambio de base. Queda, pues,
\marginpar{155}
probar que, si $F$ y $G$ son $\cal{E}'$-funtores
$\cal{F}'\to\cal{G}'$, entonces \emph{la aplicación}
$$
u\mto\pr_1\circ u
$$
\emph{induce una biyección}
$$
\Hom_{\cal{E}'}(F,G)\isomto\Hom_{\cal{E}}(\pr_1\circ F,\pr_1\circ
G)\quoi.
$$
La comprobación de este hecho es inmediata y se deja al lector.

De este isomorfismo~\eqref{eq:VI.3.i} y del final del número anterior se deduce
que
$$
\SheafHom_{\cal{E}'/-}(\cal{F}\times_{\cal{E}}\cal{E}',\cal{G}
\times_{\cal{E}}\cal{E}')
$$
\emph{puede considerarse como un funtor en}
$\cal{E}',\cal{F},\cal{G}$, \emph{de la categoría}
$\Cat_{/\cal{E}}^\circ\times\Cat_{/\cal{E}}^\circ\allowbreak\times\Cat_{/\cal{E}}$
\emph{en la categoría}~$\Cat$, isomorfo al funtor definido por la expresión
$\SheafHom_{\cal{E}/-}(\cal{F}\times_{\cal{E}}\cal{E}',\cal{G})$. En
particular, para $\cal{F}$ y~$\cal{G}$ fijos, se obtiene un funtor en~$\cal{E}'$,
\ifthenelse{\boolean{orig}}
{$\cal{E}'\to\SheafHom_{\cal{E}''/-}(\cal{F}',\cal{G}')=
\SheafHom_{\cal{E}'/-}(\cal{F}\times_{\cal{E}}
\cal{E}',\cal{G}\times_{\cal{E}}\cal{E}')$}
{$\cal{E}'\mto\SheafHom_{\cal{E}'/-}(\cal{F}',\cal{G}')=
\SheafHom_{\cal{E}'/-}(\cal{F}\times_{\cal{E}}
\cal{E}',\cal{G}\times_{\cal{E}}\cal{E}')$}, y, en particular, el
$\cal{E}$-funtor de proyección $\lambda\colon\cal{E}'\to\cal{E}$ define un
morfismo, es decir, un funtor
$$
\lambda^*_{\cal{F},\cal{G}}\colon
\SheafHom_{\cal{E}/-}(\cal{F},\cal{G})\to
\SheafHom_{\cal{E}'/-}(\cal{F}',\cal{G}')
$$
que vamos a explicitar. En los conjuntos de objetos de ambos miembros, se trata
de la aplicación
$$
f\mto f'=f\times_{\cal{E}}\cal{E}'
$$
que expresa la dependencia funtorial de
$\cal{F}\times_{\cal{E}}\cal{E}'$ respecto del objeto~$\cal{F}$
sobre~$\cal{E}$. Por otra parte, consideremos dos $\cal{E}$-funtores
$$
f,g\colon\cal{F}\to\cal{G}
$$
y un homomorfismo de $\cal{E}$-funtores
$$
u\colon f\to g\quoi,
$$
vamos
\marginpar{156}
a explicitar el homomorfismo correspondiente de $\cal{E}'$-funtores:
$$
u'\colon f'\to g'\quoi.
$$
Para todo
$$
\xi'=(\xi,S')\in\Ob(\cal{F}')
$$
con
$$
\xi\in\Ob(\cal{F}),\quad S'\in\Ob(\cal{E}'),\quad p(\xi)=\lambda(S')=S
$$
el morfismo
$$
u'(\xi')\colon f'(\xi')=(f(\xi),S')\to g'(\xi')=(g(\xi),S')\quad
\text{en }\cal{G}'
$$
se define mediante la fórmula
$$
u'(\xi')=(u(\xi),\id_{S'})
$$
(lo cual es, en efecto, un $S'$-morfismo en~$\cal{G}'$, pues
$q(u(\xi))=\lambda(\id_{S'})=\id_S$).

Consideremos ahora un $\cal{E}$-funtor cualquiera
$$
\lambda'\colon\cal{E}''\to\cal{E}'
$$
y el funtor correspondiente
$$
\SheafHom_{\cal{E}'/-}(\cal{F}\times_{\cal{E}}\cal{E}',
\cal{G}\times_{\cal{E}}\cal{E}')\to
\SheafHom_{\cal{E}''/-}(\cal{F}\times_{\cal{E}}\cal{E}'',
\cal{G}\times_{\cal{E}}\cal{E}'')\quoi,
$$
Afirmamos que este funtor no es sino el obtenido por el procedimiento anterior,
partiendo de~$\cal{F}'$ y $\cal{G}'$ sobre~$\cal{E}'$ y considerando
$\cal{E}''$ como una categoría sobre~$\cal{E}'$, habida cuenta de los
isomorfismos de \emph{«transitividad del cambio de base»}:
$$
\ifthenelse{\boolean{orig}}
{\cal{F}'\times_{\cal{E}'}\cal{E}''\isomto \cal{F}''=
\cal{F}\times_{\cal{E}}\cal{E}''\quad
\text{y}\quad\cal{G}'\times_{\cal{E}}\cal{E}''\isomto\cal{G}''
=\cal{G}\times_{\cal{E}}\cal{E}''\quoi,}
{\cal{F}'\times_{\cal{E}'}\cal{E}''\isomto
\cal{F}''=\cal{F}\times_{\cal{E}}\cal{E}''\quad
\text{y}\quad\cal{G}'\times_{\cal{E}'}\cal{E}''\isomto\cal{G}''
=\cal G\times_{\cal E}\cal E''\quoi,}
$$
que implican un isomorfismo canónico
$$
\SheafHom_{\cal{E}''/-}(\cal{F}'\times_{\cal{E}'}\cal{E}'',
\cal{G}'\times_{\cal{E}'}\cal{E}'')\isomto
\SheafHom_{\cal{E}''/-}(\cal{F}\times_{\cal{E}}\cal{E}'',
\cal{G}\times_{\cal{E}}\cal{E}'')
$$
La
\marginpar{157}
comprobación de esta compatibilidad es inmediata y se deja al lector.

Los funtores que se acaban de definir son compatibles con las operaciones de
composición definidas en el número anterior, en el sentido preciso siguiente:
si $\cal{F}$, $\cal{G}$ y~$\cal{H}$ son categorías sobre~$\cal{E}$ y se pone
$$
\cal{F}'=\cal{F}\times_{\cal{E}}\cal{E}'\quoi,\quad \cal{G}'=
\cal{G}\times_{\cal{E}}\cal{E}'\quoi,\quad \cal{H}'=
\cal{H}\times_{\cal{E}}\cal{E}'\quoi,
$$
el siguiente diagrama de funtores es conmutativo:
\ifthenelse{\boolean{orig}}
{\def\VIthreebase{\cal{E}}}
{\def\VIthreebase{\cal{E}'}}
$$
\xymatrix{ \SheafHom_{\cal{E}/-}(\cal{F},\cal{G})\times
\SheafHom_{\cal{E}/-}(\cal{G},\cal{H}) \ar[r]
\ar[d]_{\lambda^*_{\cal{F},\cal{G}}\times \lambda^*_{\cal{G},\cal{H}}}
& \SheafHom_{\cal{E}/-}(\cal{F},\cal{H})
\ar[d]^{\lambda^*_{\cal{F},\cal{H}}} \\
\SheafHom_{\VIthreebase/-}(\cal{F}',\cal{G}')\times
\SheafHom_{\cal{E}'/-}(\cal{G}',\cal{H}')\ar[r] &
\SheafHom_{\cal{E}'/-}(\cal{F}',\cal{H}')}
$$
donde las flechas horizontales son los funtores de composición definidos en el
número anterior. Esta conmutatividad se expresa mediante las fórmulas
$$
(gf)'=g'f'
$$
para $f\in\Hom_{\cal{E}}(\cal{F},\cal{G})$,
$g\in\Hom_{\cal{E}}(\cal{G},\cal{H})$ (fórmula que expresa simplemente la
funtorialidad del cambio de base), y mediante la fórmula
$$
(v*u)'=v'*u'
$$
cuando $u\colon f\to f_1$ es una flecha de
$\SheafHom_{\cal{E}/-}(\cal{F},\cal{G})$, mientras que
$v\colon g\to g_1$ es una flecha
de~$\SheafHom_{\cal{E}/-}(\cal{G},\cal{H})$. La comprobación de esta fórmula se
deduce fácilmente de las definiciones.

En lo que sigue, nos interesaremos sobre todo por
$\SheafHom_{\cal{E}}(\cal{F},\cal{G})$ (y por ciertas subcategorías notables de
esta) cuando $\cal{F}=\cal{E}$, e introducimos
\ifthenelse{\boolean{orig}}{por estas razones}{por esta razón} una notación
especial:
$$
\bf{\Gamma}(\cal{G}/\cal{E})=
\SheafHom_{\cal{E}}(\cal{E},\cal{G})\quoi,\quad
\Gamma(\cal{G}/\cal{E})=\Ob(\bf{\Gamma}(\cal{G}/\cal{E}))=
\ifthenelse{\boolean{orig}}
{\Hom_{\cal{E}}(\cal{G},\cal{E})\quoi.}
{\Hom_{\cal{E}}(\cal{E},\cal{G})\quoi.}
$$
\label{indnot:fm}\oldindexnot{$\bf{\Gamma}(\cal{G}/\cal{E})$ y $\Gamma(\cal{G}/\cal{E})$|hyperpage}%

\begin{remarquesstar}
Cuando
\marginpar{158}
$\cal{E}$ es una categoría puntual, es decir, cuando $\Ob(\cal{E})$ y
$\Fl(\cal{E})$ se reducen cada uno a un solo elemento, lo que significa también
que $\cal{E}$ es un objeto final de la categoría~$\Cat$, el dato de una
categoría sobre~$\cal{E}$ equivale al dato de una categoría sin más (pues hay
un único funtor de~$\cal{F}$ en~$\cal{E}$). Más precisamente,
$\Cat_{/\cal{E}}$ es entonces isomorfa a~$\Cat$. Además, las categorías
$\SheafHom_{\cal{E}/-}(\cal{F},\cal{G})$ no son sino las categorías
$\SheafHom(\cal{F},\cal{G})$. Recordemos entonces que la fórmula fundamental
$$
\Hom(\cal{H},\SheafHom(\cal{F},\cal{G}))\isomto
\Hom(\cal{F}\times\cal{H},\cal{G})
$$
(isomorfismo funtorial en los tres argumentos que figuran en ella) permite
interpretar $\SheafHom(\cal{F},\cal{G})$ axiomáticamente, en términos internos
a la categoría~$\Cat$, de manera que el formalismo conocido de las categorías
$\SheafHom$ aparece como un caso particular de un formalismo válido en
categorías tales como~$\Cat$, en las que existen «objetos $\SheafHom$ internos»
(definidos por la fórmula anterior). Hay una interpretación análoga de
$\SheafHom_{\cal{E}/-}(\cal{F},\cal{G})$ cuando se supone de nuevo que
$\cal{E}$ es arbitraria, mediante la fórmula
$$
\Hom(\cal{H},\SheafHom_{\cal{E}/-}(\cal{F},\cal{G}))\isomto
\Hom_{\cal{E}}(\cal{F}\times\cal{H},\cal{G})
$$
(isomorfismo funtorial en los tres argumentos). De este modo, las propiedades
formales expuestas en los \No \Ref{VI.2},~\Ref{VI.3} son casos particulares de
resultados más generales, válidos en las categorías en las que existen los
objetos internos $\SheafHom_{\cal{E}/-}(\cal{F},\cal{G})$ (cuando $\cal{F}$ y~$\cal{G}$
son dos objetos de la categoría de objetos sobre un tercero~$\cal{E}$).
\end{remarquesstar}
